\documentclass[10pt,twocolumn]{article}
\usepackage[T1]{fontenc}
\usepackage[utf8]{inputenc}
\usepackage[margin=0.75in,columnsep=0.25in]{geometry}
\usepackage{amsmath,amssymb}
\usepackage[demo]{graphicx}
\usepackage{booktabs}
\usepackage{cite}

\title{Adaptive Line Breaking for Incremental Document Preview}
\author{A.~Author \and B.~Coauthor\\
  Department of Typesetting, Example University\\
  \texttt{\{a.author,b.coauthor\}@example.edu}}
\date{}

\begin{document}
\maketitle

\begin{abstract}
We describe an incremental typesetting pipeline that retypesets only the
paragraph touched by a keystroke while keeping page breaks stable. On a
corpus of homework, lecture notes and short articles the median
keystroke-to-paint latency is below one frame. We report the measured
distribution, the failure modes, and an honest inventory of the constructs
the pipeline does not yet support.
\end{abstract}

\section{Introduction}
Live preview of \LaTeX{} documents is usually approximated by re-running a
batch compiler on every save. This paper takes the opposite position: the
document is a sequence of blocks, and a keystroke invalidates exactly one
block~\cite{knuth84,plass81}. The remaining blocks keep their shaped
glyph runs and their line breaks.

The contribution is threefold. First, a block cache keyed by the source
bytes of each paragraph (Section~\ref{sec:cache}). Second, a page builder
that replays TeX's penalty and glue model so that the cached blocks fall on
the same pages the batch compiler would choose (Section~\ref{sec:pages}).
Third, a measurement harness that compares every page against pdfLaTeX at
144~dpi (Section~\ref{sec:eval}).

\section{Block cache}\label{sec:cache}
A paragraph is identified by the SHA-256 of its source bytes together with
the font selection in force at its first character. Two paragraphs with
identical bytes but different fonts must not share a cache entry; the
initial version of the cache made exactly this mistake, and a bold heading
was painted in the body font for one frame.

\begin{figure}[t]
  \centering
  \includegraphics[width=\columnwidth]{block-cache}
  \caption{Cache hit rate against edit position. Edits inside display math
  invalidate the surrounding paragraph as well.}
  \label{fig:cache}
\end{figure}

Figure~\ref{fig:cache} shows the hit rate across an editing session. The
cost of a miss is dominated by shaping; the Knuth--Plass pass itself is
linear in the number of candidate breakpoints for the \emph{total-fit}
variant used here.

\section{Page builder}\label{sec:pages}
The page builder implements the contribution list of \emph{The
\TeX book}, Chapter~15. For each block we record its natural height $h$,
depth $d$ and the penalty $p$ following it. A page break at a block with
badness $b$ has cost
\begin{equation}
  c = \begin{cases}
    b + p + q & \text{if } b < \infty \text{ and } p \le -10000,\\
    \infty & \text{otherwise},
  \end{cases}
  \label{eq:cost}
\end{equation}
where $q$ is the insertion penalty. Equation~\eqref{eq:cost} is evaluated
once per contributed block, which is sufficient because the page goal is
fixed for a given class and geometry.

\subsection{Stability under edits}
Table~\ref{tab:stability} lists how often a single keystroke moved a page
break in the three corpora. Page breaks moved only when the edited
paragraph crossed a page boundary or when a float was displaced.

\begin{table}[t]
  \centering
  \caption{Page-break stability per keystroke.}
  \label{tab:stability}
  \begin{tabular}{lrrr}
    \toprule
    Corpus & Keystrokes & Moved & Rate \\
    \midrule
    Homework      & 12\,480 & 41  & 0.33\% \\
    Lecture notes & 30\,112 & 208 & 0.69\% \\
    Articles      & 8\,930  & 97  & 1.09\% \\
    \bottomrule
  \end{tabular}
\end{table}

\section{Evaluation}\label{sec:eval}
Each document in the corpus has a reference PDF produced by pdfLaTeX with
\texttt{SOURCE\_DATE\_EPOCH=0}. Both PDFs are rasterised at 144~dpi and
compared pixel by pixel. We deliberately report raw differing-pixel counts
rather than a similarity score: a score hides whether a difference is a
missing glyph or a one-pixel shift of an entire column.

\subsection{Threats to validity}
The corpus is small and written by the authors. Two-column documents
exercise \texttt{\textbackslash columnwidth} and float placement in ways a
single-column homework does not, and we make no claim about documents that
use \texttt{multicol}, \texttt{longtable} or TikZ.

\section{Conclusion}
Incremental typesetting is practical when the cache key includes the font
state and the page builder replays TeX's cost model exactly. The remaining
gap is not latency but coverage, and the coverage inventory is published
with every run.

\begin{thebibliography}{9}
\bibitem{knuth84} D.~E. Knuth, \emph{The \TeX book}. Addison-Wesley, 1984.
\bibitem{plass81} D.~E. Knuth and M.~F. Plass, ``Breaking paragraphs into
  lines,'' \emph{Software: Practice and Experience}, vol.~11, no.~11,
  pp.~1119--1184, 1981.
\end{thebibliography}

\end{document}
